% Source-derived chapter 13: End of the proof
% Original source: standard-conjectures-abelian-fourfolds
\chapter{End of the proof}
\label{standard-conjectures-abelian-fourfolds-chapter-13}
\source{standard-conjectures-abelian-fourfolds}

\begin{remark}[Source section]
\label{standard-conjectures-abelian-fourfolds-chapter-13-source}
\source{standard-conjectures-abelian-fourfolds}
\source{standard-conjectures-abelian-fourfolds}
This chapter reproduces the corresponding section of the retrieved
LaTeX source, including its definitions, statements, examples, and
proofs. It has not yet been translated into Lean.
\notready
\end{remark}

% The source section title was: End of the proof
\label{endpf}
\source{standard-conjectures-abelian-fourfolds}
We are now ready to show Theorem \ref{mainthm}  (via Proposition \ref{padicreduction}).  

\begin{proof}[Proof of Theorem \ref{mainthm}]
Thank to Proposition \ref{ordinarycase}, we can work under Assumption \ref{ipos}, and fix the positive integer $i$ as in the assumption as well as the field $F$ from Proposition \ref{propendo}.

Let us first suppose that  $F$ is unramified. We can then combine Proposition \ref{warmbrinon} and Corollary \ref{brinon} to conclude that the quadratic forms  $q_{B,p}$ and $q_{Z,p}$ are isomorphic if and only if $1/p^i$ is a norm. Because of Lemma \ref{notnorm}, this is the case   if and only if $i$ is even.


Suppose now that $F$ is ramified but it is not $\mathbb{Q}_2(\sqrt{3})$ nor $\mathbb{Q}_2(\sqrt{-1})$. Then, following  Remark \ref{remtame}, we can combine Proposition \ref{warmbrinonram} and Corollary \ref{brinonram} to conclude that the quadratic forms $q_{B,p}$ and $q_{Z,p}$ are isomorphic if and only if    $ a^i$ is a norm. Because of Lemma \ref{notnormram}, this is the case   if and only if $i$ is even.

If $F$ is  $\mathbb{Q}_2(\sqrt{-1})$ we can  combine Proposition \ref{warmbrinonram} and Corollary \ref{brinonwild} to conclude that the quadratic forms $q_{B,p}$ and $q_{Z,p}$ are isomorphic if and only if    $  (\alpha  \cdot \sigma(\alpha) )^i$ is a norm. Because of Lemma \ref{notnormwild}, this is the case   if and only if $i$ is even.

Finally, if $F$ is  $\mathbb{Q}_2(\sqrt{3})$ we can  combine Proposition \ref{warmbrinonram} and Corollary \ref{brinonwild2} to conclude that the quadratic forms  $q_{B,p}$ and $q_{Z,p}$ are isomorphic if and only if    $  (\alpha  \cdot \sigma(\alpha) )^i$ is a norm. Because of Lemma \ref{notnormwild2}, this is the case   if and only if $i$ is even.


In conclusion, we showed for that, for any possible $F$, the quadratic forms $q_{B,p}$ and $q_{Z,p}$ are isomorphic if and only if $i$ is even. This shows Theorem \ref{mainthm}  via Proposition \ref{padicreduction}.
\end{proof}


\begin{lemma}
\source{standard-conjectures-abelian-fourfolds}
\notready\label{notnorm}
\source{standard-conjectures-abelian-fourfolds}
The element $1/p\in \Q_p^*$ does not belong to the group of norms $ N_{\Q_{p^2}/\Q_p}(\Q^*_{p^2}).$
\end{lemma}
\begin{proof}
For $f\in \Q_{p^2}$,   the $p$-adic valuation of $f$ and $\varphi (f)$ coincide, hence a norm must have even $p$-adic valuation.
\end{proof}

\begin{lemma}
\source{standard-conjectures-abelian-fourfolds}
\notready\label{notnormram}
\source{standard-conjectures-abelian-fourfolds}
Suppose that $F$ is ramified but it is not $\mathbb{Q}_2(\sqrt{3})$ nor  $\mathbb{Q}_2(\sqrt{-1})$. Let $\sqrt{a}\in \Q_{p^2}$ be as in Corollary \ref{brinonram}. Then the element $a \in \Q_p^*$ does not belong to the group of norms $ N_{F/\Q_p}(F^*).$
\end{lemma}
\begin{proof}
This is about finding nonzero solutions $ x,y \in \Q_p$ of the equation
\begin{equation}\label{calcare} x^2 -\pi^2 y ^2 = a. \end{equation}
\source{standard-conjectures-abelian-fourfolds}
We can suppose that    $a$ has $p$-adic valuation zero. Then $x$ has valuation zero and $y$ has non-negative valuation. 

Consider first the case $p\neq 2$. By reducing modulo $p$ we see that the existence of a solution $(x,y)$ would imply that $a$ would be a square in $\mathbb{F}_p$  which gives a contradiction. 

Consider now the case $p= 2$. We can then take $a=5$ and $\pi^2=2,6,10$ or $14$, see Remark \ref{remtame}. If we reduce modulo $8$ the  equation (\ref{calcare}) we see that the existence of a solution $(x,y)$ would imply that  $3,5$ or $7$ would be a square in $\mathbb{Z}/ 8 \mathbb{Z}$ which gives a contradiction. 
\end{proof}





\begin{lemma}
\source{standard-conjectures-abelian-fourfolds}
\notready\label{notnormwild}
\source{standard-conjectures-abelian-fourfolds}
Let $\alpha \in \Q_{2^4}(\sqrt{3})=\Q_{2^4}(\sqrt{-1})$ be as in Corollary \ref{brinonwild}. Then the element $\alpha  \cdot \sigma(\alpha) \in \Q_2^*$ does not belong to the group of norms $ N_{\mathbb{Q}_2(\sqrt{-1})/\Q_2}(\mathbb{Q}_2(\sqrt{-1})^*).$
\end{lemma}
\begin{proof}
 Write $\alpha= \alpha_1+\sqrt{-1}\cdot\alpha_2$ with $\alpha_1,\alpha_2 \in \Q_{2^4}$, then we have
 \[\varphi(\alpha_1)=-\alpha_2,  \hspace{1cm} \varphi(\alpha_2)=\alpha_1, \hspace{1cm} \alpha  \cdot \sigma(\alpha)=\alpha_1^2+\varphi(\alpha_1^2).\]
 Let us make explicit computations. Fix the presentations
\[\mathbb{F}_{2^2}=\mathbb{F}_{2}[x]/(x^2+x+1) \hspace{1cm} \mathbb{F}_{2^4}=\mathbb{F}_{2^2}[y]/(y^2+xy+1)\]
and write
\[\mathbb{Q}_{2^2}=\mathbb{Q}_{2}(\frac{-1+\sqrt{-3}}{2}) \hspace{1cm} \mathbb{Q}_{2^4}=\mathbb{Q}_{2^2}[y]/(y^2+(\frac{-1+\sqrt{-3}}{2})y+1).\]
Then $\Delta=\sqrt{(\frac{-1+\sqrt{-3}}{2})^2-4}=\sqrt{\frac{-1-\sqrt{-3}}{2}-4}$ can be chosen as $\alpha_1$. Thus, 
\[\alpha  \cdot \sigma(\alpha)=\Delta^2 + \varphi(\Delta^2)= -9.\]
This is a norm if and only if one can find a nonzero solution $u,v \in \Q_2$ of the equation $u^2+v^2=-9$. This is impossible by looking at the equation modulo $8$.
\end{proof}

\begin{lemma}
\source{standard-conjectures-abelian-fourfolds}
\notready\label{notnormwild2}
\source{standard-conjectures-abelian-fourfolds}
Let $\alpha \in \Q_{2^4}(\sqrt{3})=\Q_{2^4}(\sqrt{-1})$ be as in Corollary \ref{brinonwild2}. Then the element $\alpha  \cdot \sigma(\alpha) \in \Q_2^*$ does not belong to the group of norms $ N_{\mathbb{Q}_2(\sqrt{3})/\Q_2}(\mathbb{Q}_2(\sqrt{3})^*).$
\end{lemma}
\begin{proof}
 We argue as in the proof of Lemma \ref{notnormwild}. Write $\alpha= \alpha_1+\sqrt{-1}\cdot \alpha_2$ with $\alpha_1,\alpha_2 \in \Q_{2^4}$, then we have
\[\varphi(\alpha_1)=-\alpha_2,  \hspace{1cm} \varphi(\alpha_2)=\alpha_1, \hspace{1cm}\]

 Moreover, as the elements fixed by $\sigma$ are exactly those in $B_{\crys}$, we have
 \[ \sigma(\sqrt{-1})=-\sqrt{-1} \hspace{1cm} \sigma(\alpha_1)= \alpha_1 \hspace{1cm} \sigma(\alpha_1)= \alpha_1 \]
 and we deduce that $\alpha  \cdot \sigma(\alpha)=\alpha_1^2+\varphi(\alpha_1^2).$ This means that we can use the same computations as in  the proof of Lemma \ref{notnormwild} and for the same choice of $\alpha$ in there we have \[\alpha  \cdot \sigma(\alpha) = -9.\]
This is a norm if and only if one can find a nonzero solution $u,v \in \Q_2$ of the equation $u^2 - 3 v^2=-9$. This is impossible by looking at the equation modulo $8$.
\end{proof}

\begin{remark}
\source{standard-conjectures-abelian-fourfolds}
\notready
Let us point out that the formulation of Proposition \ref{padicreduction} does not involve the ring $B_{\crys}$. We do not know if there is a way to show Theorem \ref{mainthm} via   this proposition without computing explicitly the $p$-adic periods.
\end{remark}


\begin{appendices}
